% !TEX root =paper.tex


\section{Matching}
\label{sec:matching}

%\Anna{Constants: $\eps_{F}<1/20, \eps_B\geq 3\eps_{1/2}, \alpha\geq 3\eps_\eta$}
%\Shayan{MOVE the FACT}
%\begin{fact}
%\label{fact:ErightarrowS}
%Let $S$ be a $2 + \eps_\eta$ near min-cut. Then $x(E^\rightarrow (S)) \ge |S|- 1 - \eps_\eta/2$.
%\end{fact}
%\begin{proof}
%$$
%x(E^\rightarrow (S))  = \frac12 \sum_{a \in S} (x(\delta (a)) - x(\delta ^\uparrow (a))) 
%\ge \frac{2 |S| - (2 + \eps_\eta)}{2}
%= |S|- 1 - \eps_\eta/2.$$
%\end{proof}

%\Nathan{Can this be stated more generally?}

The main result of this section is to construct a matching that we use in order to decide which edges will have positive slack to compensate for the negative slack of edges going higher. 
Refer to \cref{ex:simple} for a high-level motivation to construct a matching. . 

\begin{definition} [$\eps_F$ fractional edge]
\hypertarget{tar:epsF}{For $z\geq 0$ we say that $z$ is $\eps_{F}$-fractional if $\eps_F \le z\leq 1-\eps_{F}$.}
\end{definition}
\noindent The following lemma is the main result of this section
\begin{lemma}[Matching Lemma]
\label{lem:matching}
For any $S \in \cH$, $\eps_{F}\leq 1/10, \eps_B\geq 21\eps_{1/2}, \alpha\geq 2\eps_\eta$, $\eps_{1/2}\leq 0.0002$, there is a matching from good edges (see \cref{def:goodedges}) in $E^\rightarrow(S)$ to edges in $\delta(S)$ where every good edge bundle $\bbe={\bf(u,v)}$ (where $u,v\in\cA(S)$) is matched to a fraction $m_{\bbe,u}$ of edges in $\delta^\uparrow(u)$ and a fraction $m_{\bbe,v}$ of $\delta^\uparrow(v)$, and:
\begin{eqnarray}
\label{eq:uvbadematch}m_{\bbe,u} F_u  + m_{\bbe,v} F_v &\le & x_\bbe (1 + \alpha)\\
\label{eq:matchedamount}\sum_{e \in \delta^\rightarrow (u)} m_{\bbe,u} &=& x(\delta^\uparrow(u)) Z_u,
\end{eqnarray}
where for every atom $u\in \cA(S)$, define
$$\hypertarget{tar:Fu}{F_u=1-\eps_B\I {x(\delta^\uparrow(u))\text{ is $\eps_F$ fractional}},} \quad
\hypertarget{tar:Zu}{Z_u:=\left(1 + \I{|\cA(S)|\geq 4, x(\delta^\uparrow(u))\leq \eps_F}\right).}$$
%which is active when $\delta^\uparrow(u))$ is very close to {\em zero}.
\end{lemma}








Roughly speaking, the intention of the above lemma is to match good edges in $E^\rightarrow(S)$ to a similar fraction of edges that go higher (such that an edge bundle $e$ adjacent to atoms $u,v$ is only matched to edges in $\delta^\uparrow(u),\delta^\uparrow(v)$). Since we never ``reduce'' bad edges in the proof of payment theorem (\cref{thm:payment-main}), we don't use them in the matching. That inherently can cause a problem, as there could not be ``enough'' good edges in $E^\rightarrow(S)$ to saturate the edges going higher in the matching. The parameter $F_u$ help us in this regard; in particular, it allows us to match some of the (good) edges in $E^\rightarrow(S)$ to more than their fraction in $\delta(S)$.
 
Next, we motivate the parameter $Z_u$. If  $x(\delta^\uparrow(u))\approx 0$, when those edges are reduced the conditional probability that $\delta(u)_T$ is even could be very close to 0. 
The parameter $Z_u$ lets us match twice as many edges to $\delta^\uparrow(u)$; so there will be only half a burden to fix the parity of $\delta(u)_T$.
See the discussion in \hyperlink{xu-close-to-zero}{overview section} for more details. 

%As we will see in \cref{sec:payment}, the downside that we need a stronger conditional 

% \cref{eq:uvbadematch} says that if $x(\delta^\uparrow(u))$ is fractional then edges incident to $u$ can be matched to a larger faction of edges in $\delta^\uparrow(u)$. 

Throughout this section we adopt the following notation: For a cut $S\in \cH$ and a set $W\subseteq \cA(S)$, we write 
\begin{align*}
E(W,S\smallsetminus W)&:=\cup_{u\in W, v\in \cA(S)\smallsetminus W} E(u,v),\\
\delta^\uparrow(W)&:=\cup_{u\in W} \delta^\uparrow(u)=\delta(W)\cap \delta(S), \\
\delta^\rightarrow(W) &:= \cup_{u\in W} \delta^\rightarrow(u).
\end{align*}
Note that in $\delta^\rightarrow(W)\not\subseteq \delta(W)$ since it includes edge bundles between atoms in $W$. 

Before proving the main lemma we record the following facts.
\begin{lemma}
\label{fact:anothercutineq}
For any $S\in \cH$ %with children $u_1,\dots,u_k$ 
and $W \subsetneq \cA(S)$ (recall $\cA(S)$ is the set of $u \in \cH$ with $\p(u) = S$), we have 
$$ x(\delta^\rightarrow(W))  %+ x(E(W, S \smallsetminus W) 
	\geq \frac12 \sum_{u\in W} x(\delta(a)) - \eps/2 \ge |W| - \eps/2.$$
	%(Recall also that $E^\rightarrow (W)$ between atoms in $W$.)
\end{lemma}
\begin{proof}
We have
$$ x(\delta^\rightarrow (W))  % + x(E(W, S \smallsetminus W)  
= \frac12 \left(\sum_{u \in W} (x(\delta (u)) +   x(	E(W, S\smallsetminus W)) -  x(\delta ^\uparrow (W))\right).$$
Since $x(\delta (S\smallsetminus W)) \ge 2$ and $x(\delta (S)) \le 2 + \eps$, we have:
$$\text{(a) }  x(	E(W, S\smallsetminus W)) + x(\delta^\uparrow(S\smallsetminus W)))\ge 2 \quad\text{ and }\quad\text{(b) }  x(\delta^\uparrow(W)) +  x(	\delta^\uparrow( S\smallsetminus W)) \le 2 + \eps.$$
Subtracting (b) from (a), we get
$$x(	E(W, S\smallsetminus W)) - x(\delta ^\uparrow (W)) \ge -\eps,$$
which after substituting into the above equation,
completes the proof of the first inequality in the lemma statement. The second inequality follows from the fact that $\delta (u) \ge 2$ for each atom $u$.
\end{proof}


\begin{lemma}\label{lem:three-nodes-good}
For $S \in\cH$, if $|\cA(S)|=3$ then there are no bad edges in \hyperlink{tar:ErightarrowS}{$E^\rightarrow(S)$}.
\end{lemma}
\begin{proof}
Suppose  $\cA(S)=\{u,v,w\}$ and
$\bbe=(u,v)$  is a bad edge bundle. Then $|x_\bbe - \frac{1}{2}| \le \epsilon_{1/2}$. In addition, by \cref{thm:badedges}, $x(\delta^\uparrow(u)),x(\delta^\uparrow(v)) \le 1/2+9\epsilon_{1/2}$. Therefore, 
$$x_{\bf(u,w)} = x(\delta(u)) - x_{\bbe} -  x(\delta^\uparrow(u))\geq 1-10\epsilon_{1/2}.$$
Similarly, $x_{\bf(v,w)}\geq 1-10\eps_{1/2}$.
Finally, since $x(\delta (S)) \ge 2$, and $x(\delta^\uparrow(u)),x(\delta^\uparrow(v))\leq 1/2+9\eps_{1/2}$, we must have $x(\delta^\uparrow(w)) \ge 1-18\epsilon_{1/2}$.
But, this contradicts the assumption that $w\in \cH$ must satisfy $x(\delta(w))\leq 2+\epsilon_\eta$.
\end{proof}

%\begin{fact}
%In any tree and for any vertices $u,v,w$, it cannot be that both $u$ is on the $v-w$ path and $w$ is on the $u-v$ path.	
%\end{fact}
%\begin{proof}
%The path $v-w$ includes $u$. But then the $u-v$ path is a subpath of the $v-w$ path, so it cannot include $w$: contradiction. 
%\end{proof}

%Let $\epsilon_\eta$ be the constant matching the largest near minimum cut in our laminar family 
%We expect $\epsilon_\eta = 4\eta$. Let $\epsilon_{1/2}$ be the edge bounding edges away from 1/2. 

%\begin{lemma}
%For a degree cut $S$, there is a matching from good edges in $E(S)$ to edges in $\delta(S)$ such that every edge of fraction $x_e$ is matched to at most fraction $x_e(1+6\epsilon_{1/2})$ when $\epsilon_{1/2} \le 1/10$ and $\epsilon_{1/2} \ge \epsilon_{\eta}$. 
%\end{lemma}
%\begin{proof}
%	We will prove this by setting up a max-flow min-cut problem. Construct a graph with vertex set $\{s,X,Y,t\}$, where $s,t$ are the source and sink. Let $X$ be the set of good edges in $E^\rightarrow(S)$ and $Y$ be the set of edges in $\delta(S)$. For every edge $e \in X$, add an arc from $s$ to $e$ of capacity $c(s,e):= (1+\alpha)x_e$ (for some $\alpha>0$ that we choose later).  Connect every $f \in Y$ to $t$ with an arc of capacity $c(f,t):=x_f$. 
%Finally, connect $e\in X$ to $f\in Y$ with a directed edge of infinite capacity, i.e., $c(e,f):=\infty$, if $f\in\delta^\uparrow(e)$.
%Now, our goal is to show that there is a flow saturating $t$, i.e. there is a flow of value $\sum_{e \in Y} x_e\leq2+\eps_\eta$.  
%
%Fix a $s$-$t$ cut $A,\overline{A}$ where $s\in A, t\in \overline{A}$. 
% We need to show $\text{cap}(A,\overline{A})\geq 2+\epsilon_\eta$.
%
%First we prove this for $A=\{s\}$. By \cref{lem:one-of-two-good} there is at most one bad 1/2 edge adjacent to every vertex. Therefore there is at most  $|S|/2$ of bad edges (the bound is met if they form a perfect matching) which contributes to at most a total of $(1/2+\eps_{1/2})|S|/2=:x_B$ fraction. If $|S|=3$ then by \cref{lem:three-nodes-good} all edges are good. Then there is at least a fraction $2-\epsilon_\eta/2$ to match to fraction $x(\delta(S))\leq 2+\eps_\eta$, and  for $\alpha\geq \eps_\eta$ we have
%$$(2-\epsilon_\eta/2)(1+\alpha) \ge 2+\epsilon_\eta$$
%as desired. 
%Otherwise there is a fraction of at least $x_G:=|S|-1-\epsilon_\eta/2-x_B$ of good edges and $|S|\ge 4$. But then we will choose $\alpha$ such that 
%$$(1+\alpha)x_G  = (1+\alpha)\left(|S| (\frac34 - \frac{\eps_{1/2}}{2}) -1 -\frac{\eps_{\eta}}{2}\right)\ge 2+\epsilon_\eta.$$
%Note that $\alpha \geq \frac{\eps_{1/2}}{2} + 2\eps_{\eta}$ clearly suffices.
%
%From now on, we assume that $A\neq \{s\}$. In the following we will prove that for any set of \textit{atoms} $T \subsetneq S$, we have:
%\begin{equation}\label{eq:goalTdeltarigthup}
%c(s,E^\rightarrow(T)) + c(s,\delta^\rightarrow(T))  = (1+\alpha)(x_G(E^\rightarrow(T)) + x_G(\delta^\rightarrow(T))) \geq x(\delta^\uparrow(T)) 
%\end{equation}
%where for a set $F$ of edges we write $x_G(F)$ to denote the sum of fraction of good edges in $F$.
%
%%We previously proved it for $T=S$.
%Let $A_X=A\cap X, A_Y=A\cap Y$ and so on. Assuming the above inequality, let us prove the lemma: First, for the set of edges $A_X$ chosen from $X$, we must choose all adjacent edges in $Y$ due to the infinite capacity edges, i.e., $A_Y$ contains the set of neighbors of $A_X$ in $Y$. Say edges in $A_Y$ will correspond to a set of vertices $Q$, i.e., $A_Y=\delta^\uparrow(Q)$. 
%%Without loss of generality we also include every $u\in S$ where $\delta^\uparrow(u)=\emptyset$ in $Q$. This, would imply that, in the worst case, $A_X=E(Q)$. 
%It follows that $A_X\subseteq E^\rightarrow(Q)$.
%Let $T =  S \smallsetminus Q$. Note that $T\neq S$. Then:
%\begin{eqnarray*}
% \text{cap}(A,\overline{A}) &= & c(A_Y,t) + c(s,\overline{A}_X)\\
% &\geq & x(\delta^\uparrow(Q)) + c(s,E^\rightarrow(T)) + c(s,E(Q,T)) \\
% &=& x(\delta^\uparrow(S))-x(\delta^\uparrow(T)) + c(s,E^\rightarrow(T)) + c(s,\delta^\rightarrow(T)) \geq x(\delta^\uparrow(S)),
%\end{eqnarray*}
%where the last inequality follows by \eqref{eq:goalTdeltarigthup}.
%%We obtain the equality due to the fact that the set of edges corresponding to the vertex set $T$ may not have any edges from the cut, as these would each contribute infinity to the capacity. Therefore we know that the set $T$ must include all possible edges which is the union of $\delta(T)$ and $E(T)$. 
%
%In the rest of the proof we show \eqref{eq:goalTdeltarigthup}.
%First, we prove the following claim.
%\begin{claim} For any $T\subsetneq S$ we can write,
%\begin{equation}\label{eq:cut-ineq}
%	x(\delta^\uparrow(T)) + x(E^\rightarrow(T))  
%	\leq \frac12 \sum_{a\in T} x(\delta(a)) + \eps_\eta/2.
%	%\le |T|(1+\epsilon_\eta/2) + \epsilon_\eta/2.
%\end{equation}
%where the sum is over all atoms of $T$.
%\end{claim}
%\begin{proof}
%%For any set we have fraction $\sum_{a \in T} x(\delta(v)) - x(E^\rightarrow(T)) - x(\delta^\uparrow(T))$ staying below.  
%First, observe
%\begin{eqnarray*}
%	x(\delta(S \smallsetminus T)) &=& x(\delta(S))  - x(\delta^\uparrow(T))+ x(\delta^\rightarrow(T)) \\
%	&\le& 2+\epsilon_\eta - x(\delta^\uparrow(T)) + \sum_{a \in T} x(\delta(a)) - 2x(E^\rightarrow(T)) - x(\delta^\uparrow(T))\\
%	&=& 2+\eps_\eta - 2x(\delta^\uparrow(T)) - 2x(E^\rightarrow(T)) + \sum_{a\in T} x(\delta(a))
%\end{eqnarray*}
%However since $T\neq S$, $x(\delta(S \smallsetminus T)) \ge 2$.
%Dividing both sides by 2 %and using $x(\delta(a))\leq 2+\eps_\eta$ for all atoms in $S$ 
%proves the claim. 
%%$$\sum_{a \in T} x(\delta(a)) + \epsilon_\eta - 2x(\delta^\uparrow(T)) - 2x(E^\rightarrow(T)) \ge 0$$
%%And rearranging 
%%$$\sum_{a \in T} x(\delta(a)) - x(E^\rightarrow(T)) - x(\delta^\uparrow(T)) \ge x(\delta^\uparrow(T)) + x(E^\rightarrow(T)) - \epsilon_\eta$$
%%So we are done provided there are no bad edges since the capacity of the LHS is multiplied by $(1+\epsilon_\eta)$. In the remaining proof we can therefore assume that $E^\rightarrow(T),\delta^\rightarrow(T)$ contain bad edges.  
%\end{proof}
%Next, we use the above claim to prove \eqref{eq:goalTdeltarigthup}.
%Suppose atoms in $T$ are adjacent to $k$ bad edges. Each of these bad edges have fraction at most $1/2+\eps_{1/2}$. 
%\begin{align}
%x_G(E^\rightarrow(T)) + x_G(\delta^\rightarrow(T))	 &= x(E^\rightarrow(T)) + x(\delta^\rightarrow(T)) - x_B(E^\rightarrow(T)) - x_B(\delta^\rightarrow(T))\notag\\
%& =  \sum_{a\in T} x(\delta(a)) - x(\delta^\uparrow(T)) - x(E^\rightarrow(T)) -x_B(E^\rightarrow(T)) - x_B(\delta^\rightarrow(T))\notag\\
%&\geq  \frac12 \sum_{a\in T} x(\delta(a)) - \eps_\eta/2 - k (1/2+\eps_{1/2})\notag\\
%&\geq  |T| - \eps_\eta/2 - k(1/2+\eps_{1/2}) \label{eq:xGEdeltalower}
%\end{align}
%In the first inequality we used \eqref{eq:cut-ineq}, and in the second inequality we used $x(\delta(a))\geq 2$ for all atoms of $S$.
% 
% On the other hand, recall that if $e=(a,b)$ is a bad edge, then $x(\delta^\uparrow(a)),x(\delta^\uparrow(b))\leq 1/2+\eps_{1/2}$ up. Therefore, we can write,
% $$ x(\delta^\uparrow(T)) \leq k(1/2+\eps_{1/2}) + (|T|-k)(1+\eps_{\eta})$$
% Now, to prove \eqref{eq:goalTdeltarigthup}, using \eqref{eq:xGEdeltalower}, it is enough to choose $\alpha$ such that,
% $$
% (1+\alpha) (|T| - \eps_{\eta}/2 - k(1/2+\eps_{1/2}) \geq k(1/2+\eps_{1/2}) + (|T|-k)(1+\eps_{\eta}),
% $$
% or equivalently,
% $$ |T|(\alpha-\eps_{\eta}) \geq k (\alpha/2 + 2\eps_{1/2} + \alpha \eps_{1/2}+\eps_{\eta}) + \eps_{\eta}/2$$
% Note that since every atom is adjacent to at most one bad edge, $k\leq |T|$. So, the above inequality simply follows because $|T|\geq 1$ and using $\alpha\geq 6\eps_{1/2}$ and $\eps_\eta < \frac1{10}\eps_{1/2}$.
%\end{proof}
%
%
%








\begin{proof}[Proof of \cref{lem:matching}]
	We will prove this by setting up a max-flow min-cut problem. 
	%To simplify notation, let 
	%$$ D:= 1+ \alpha; \quad F_u := (1-\eps_B \I {u\in F(S)})\text{ and } Z_u := 1 + \I{|S|\geq 4, x(\delta^\uparrow(u))\leq \eps_F}.$$
	Construct a graph with vertex set $\{s,X,Y,t\}$, where $s,t$ are the source and sink. We identify $X$ with  the set of good edge bundles in $E^\rightarrow(S)$ and $Y$ with the set of atoms in $\cA(S)$. For every edge bundle $\bbe \in X$, add an arc from $s$ to $\bbe$ of capacity $c(s,\bbe):= (1+\alpha) x_\bbe$.  For every $u\in \cA(S)$, %there is a vertex 
		%$u$ in $Y$ (think of $u$ as representing the edge bundle of all edges going higher from $u$) and an 
		there is an arc $(u,t)$ with capacity
		$$ c(u,t) = x(\delta^\uparrow(u))F_u Z_u.$$
		% If $u$ is incident to a bad edge in $S$, then we set the capacity $c(f_u,t):=x(\delta^\uparrow(u))(1-\eps_B)$ and we set it to $c(f_u,t):=x(\delta^\uparrow(u))$ otherwise 
		%We will drop the sub-script $u$ from $u$ when it is not necessary.
		
Finally, connect $\bbe=(u,v)\in X$ to nodes $u$ and $v\in Y$ with a directed edge of infinite capacity, i.e., $c(\bbe,u)=c(\bbe,v)=\infty$.
We will show below that there is a flow saturating $t$, i.e. there is a flow of value 
$$c(t):=\sum_{u \in \cA(S)} c(u,t) = \sum_{u\in \cA(S)} x(\delta^\uparrow(u))F_uZ_u.   $$
%x_e\leq2+\eps_\eta$.  

Suppose that in the corresponding max-flow, there is a flow of value $f_{\bbe,u}$ on the edge $(\bbe,u)$.
Define $$m_{\bbe,u} := \frac{f_{\bbe,u}}{F_u}.$$
Then \eqref{eq:uvbadematch} follows from the fact that the flow leaving $\bbe$ is at most the capacity of the edge from $s$ to $\bbe$, and \eqref{eq:matchedamount} follows by conservation of flow on the node $u$ (after cancelling out $F_u$ from both sides).



We have left to show that for any $s$-$t$ cut $A,\overline{A}$ where $s\in A, t\in \overline{A}$ that the capacity of this cut is at least $c(t)$.

\begin{claim} If $A=\{s\}$, then capacity of $(A,\overline{A})$ is at least $c(t)$.	
\end{claim}
\begin{proof}
First, note that 
\begin{align}
	c(t) &= \sum_{u \in \cA(S)} x(\delta^\uparrow(u))F_uZ_u \le \sum_{u \in \cA(S)} x(\delta^\uparrow(u))Z_u\nonumber \\ &\le \I{|\cA(S)| \ge 4} \cdot |\{u \in \cA(S): x(\delta^\uparrow(u)) \le \eps_F\}| \cdot \eps_F + x(\delta(S))\nonumber \\
	&\le 2+\eps_\eta + \eps_F\I{|\cA(S)| \ge 4}|\cA(S)| \label{eq:bound-ct}
\end{align}
because $F_u \le 1$ and $Z_u = 1 + \I{|\cA(S)| \ge 4,x(\delta^\uparrow(u)) \le \eps_F}$. 

Second, note that
\begin{align*}x(E^\rightarrow (S))  = \frac12 \sum_{u \in \cA(S)} (x(\delta (u)) - x(\delta ^\uparrow (u))) 
\ge \frac{2 |\cA(S)| - (2 + \eps_\eta)}{2}
= |\cA(S)| - 1 - \eps_\eta/2.
\end{align*}
Therefore, if there are $k$ bad edges in $E^\rightarrow(S)$, then
\begin{align}
	x_G \ge |\cA(S)| - 1 - \eps_\eta/2 - k(\frac{1}{2}+\eps_{1/2}) \label{eq:sum-inside}
\end{align}


\textbf{Case 1: $|\cA(S)|=3$.} Then $Z_u=1$ for all $u\in\cA(S)$ and by \cref{lem:three-nodes-good} all edges are good. 
%Also,
%\begin{align*}x(E^\rightarrow (S))  = \frac12 \sum_{u \in \cA(S)} (x(\delta (u)) - x(\delta ^\uparrow (u))) 
%\ge \frac{2 |\cA(S)| - (2 + \eps_\eta)}{2}
%= 2 - \eps_\eta/2.
%\end{align*}
%%by \cref{fact:ErightarrowS},
So, by \cref{eq:sum-inside},
$x(E ^\rightarrow (S))\ge 2- \eps_\eta/2$. %, and $x(\delta(S))\leq 2+\eps_\eta$. 
Thus, for $\alpha\geq 2\eps_\eta$ we have
$$c(s) = (1+ \alpha) x_G \ge (2-\epsilon_\eta/2)(1+\alpha) \ge 2+\epsilon_\eta \underset{\cref{eq:bound-ct}}{\ge} c(t)$$
as desired. 

\textbf{Case 2: $|\cA(S)|\geq 5$.} By \cref{thm:badedges} there is at most one bad half edge adjacent to every vertex. Therefore there are at most  $|\cA(S)|/2$ bad edges, so by \cref{eq:sum-inside}, 
$$(1+\alpha)x_G \ge (1+\alpha)\left(|\cA(S)| - 1 - \eps_\eta/2 - \frac{1}{2}|\cA(S)|(\frac{1}{2}+\eps_{1/2})\right) \ge 2+\epsilon_\eta + \eps_F |\cA(S)| \underset{\cref{eq:bound-ct}}{\ge} c(t)$$
%  in total (the bound is met if they form a perfect matching) which contributes to at most a total of $(1/2+\eps_{1/2})|\cA(S)|/2=:x_B$ fraction.
%So, there is a fraction of at least $x_G:=x(E ^\rightarrow (S)) - x_B \ge |\cA(S)|-1-\epsilon_\eta/2-x_B$ of good edges. %But then we will choose $\alpha$ such that 
%If $|\cA(S)|\ge 5$  then 
%So,
%\begin{align*}
%	(1+\alpha)x_G  &\geq (1+\alpha)\left(|\cA(S)| (\frac34 - \frac{\eps_{1/2}}{2}) -1 -\frac{\eps_{\eta}}{2}\right)\\
%	& \ge (1+ \alpha)\left(|\cA(S)|(\frac{3}{4} - \frac{\eps_{1/2}}{2} - \eps_F) -1- \frac{\eps_{\eta}}{2}   + \eps_F |\cA(S)|.\right)\\
%	&\ge 2+\epsilon_\eta + \eps_F |\cA(S)| \underset{\cref{eq:bound-ct}}{\ge} c(t).
%\end{align*}
where the second to last inequality holds, using $\alpha \ge  2 \eps_\eta$, $|\cA(S)|\ge 5$, $\eps_{1/2} \le 0.01$, and $\eps_F \le 0.1$.

\textbf{Case 3: $|\cA(S)|=4$, and we have 0 or 1 bad edges.} Then by \cref{eq:sum-inside}, $x_G\geq 2.5-\eps_\eta /2 - \eps_{1/2}$, so by \cref{eq:bound-ct}, $(1+\alpha)x_G \geq 2+\eps_\eta + 4\eps_F \ge c(t)$ for $\eps_F \le 0.1$, $\alpha \ge 2\eps_\eta, \eps_{1/2} \le 0.01$. %(and noting that with $|\cA(S)|=4$ at most 2 nodes can have $\delta^\uparrow(u) \le \eps_F$).
%Finally, if $|\cA(S)|=4$ and we have two bad edges, the capacity of all edges going higher is multiplied by $(1-\eps_B)$. In this case we claim that for any $u\in S$, $x(\delta^\uparrow(u))\geq 1/10\geq \eps_F$. This is because if $x(\delta^\uparrow(u))<1/10$, then there is a node $u'\in S$ such that $x(\delta^\uparrow(u'))\geq 3/5>1/2+\eps_{1/2}$. So $u'\notin B(S)$ which is a contradiction. Therefore, in this case $\I{|\cA(S)|\geq 4, x(\delta^\uparrow(u))\leq \eps_F}=0$ for all $u\in S$. So, by Lemma ?? % if more than 1/2 + alpha goes up, then no bad edges 
%$$(1+\alpha)x_G \geq (1+ \alpha)(2-2\eps_{1/2} - \frac{\eps_\eta}{2}) \geq (2+\eps_{\eta})(1-\eps_B)$$
%for $\eps_B\ge\eps_{1/2}$.


\textbf{Case 4: $|\cA(S)|=4$, and there are 2 bad edges.}
Then they form a perfect matching inside $S$ and  for each
$u\in\cA(S)$, $x(\delta^\uparrow(u))\le 1/2 + 9\eps_{1/2}$ (see \cref{thm:badedges}).

Therefore it must also be the case that $x(\delta^\uparrow(u))\ge  \eps_F$ for each $u\in\cA(S)$. If not,  there would have to be a node $u' \in \cA(S)$ such that
$x(\delta^\uparrow(u')) \ge (2- \eps_F)/3 > 1/2 + 9\eps_{1/2}$,
which is a contradiction to $u'$ having an incident bad edge.
Thus, for each $u \in \cA(S)$, $x(\delta^\uparrow(u))$ is $\eps_F$-fractional, i.e., $F_u=1-\eps_B$ and $Z_u=1$ implying that 
 $c(t) \le (2 + \eps_\eta)(1- \eps_B)$.
 Therefore, by \cref{eq:sum-inside},
 $$c(s) = (1+ \alpha) x_G \ge (1+ \alpha) (2-2\eps_{1/2} - \eps_\eta/2,)$$ 
 and the rightmost quantity is at least $c(t)$ for $\eps_B \ge 2 \eps_{1/2}$ and $\alpha \ge 2 \eps_\eta$.
\end{proof}




From now on, we assume that the min s-t cut $A\neq \{s\}$. In the following we will prove that for any set of \textit{atoms} $W \subsetneq S$, we have:
\begin{equation}\label{eq:goalTdeltarigthup}
 c(s,\delta^\rightarrow(W))  = (1+\alpha) x_G(\delta^\rightarrow(W)) \geq c(\delta^\uparrow(W),t) 
\end{equation}
where for a set $F$ of edges we write $x_G(F)$ to denote the total fractional value of good edges in $F$.

%We previously proved it for $T=S$.
Let $A_X=A\cap X, A_Y=A\cap Y$ and so on. Assuming the above inequality, let us prove the lemma: First, for the set of edges $A_X$ chosen from $X$, let $Q$ be the set of endpoints of all edge bundles in $A_X$ (in $\cA(S)$). %So, $A_X\subseteq E^\rightarrow(Q)$.
%Without loss of generality we also include every $u\in S$ where $\delta^\uparrow(u)=\emptyset$ in $Q$.

Observe that we must choose all atoms in $Q$ inside $A_Y$ due to the infinite capacity arcs, i.e., $Q\subseteq A_Y$. %Say edges in $A_Y$ will correspond to a set of vertices $Q$, i.e., $A_Y=\delta^\uparrow(Q)$. 
% This, would imply that, in the worst case, $A_X=E^\rightarrow(Q)$. 
Let $W =  S \smallsetminus Q$. Note that $W\neq S$. Then:
\begin{eqnarray*}
 c(A,\overline{A}) &= & c(A_Y,t) + c(s,\overline{A}_X)\\
 &\geq  & c(\delta^\uparrow(Q),t) + c(s,\delta^\rightarrow(W))  \\
 &=& c(\delta^\uparrow(S),t)-c(\delta^\uparrow(W)) + c(s,\delta^\rightarrow(W)) \geq c(\delta^\uparrow(S),t),
\end{eqnarray*}
where the last inequality follows by \eqref{eq:goalTdeltarigthup}.
%We obtain the equality due to the fact that the set of edges corresponding to the vertex set $T$ may not have any edges from the cut, as these would each contribute infinity to the capacity. Therefore we know that the set $T$ must include all possible edges which is the union of $\delta(T)$ and $E(T)$. 

Finally, we prove \eqref{eq:goalTdeltarigthup}. Suppose atoms in $W$ are adjacent to $k$ bad edges. Then
\begin{align}
 x_G(\delta^\rightarrow(W))	 &=  x(\delta^\rightarrow(W))  - x_B(\delta^\rightarrow(W))\notag\\
\intertext{which by \cref{fact:anothercutineq} and the fact that each bad edge has fraction at most $1/2+\eps_{1/2}$, is}
&\geq  |W| - \eps_\eta/2 - k(1/2+\eps_{1/2}). \label{eq:xGEdeltalower}
\end{align}
%since $x(\delta (a)) \ge 2$ for each atom $a$ in $T$.
%In the first inequality we used \eqref{eq:cut-ineq}, and in the second inequality we used $x(\delta(a))\geq 2$ for all atoms of $S$.
 
 To upper bound $ c(\delta^\uparrow(W),t)$, we observe that for any $u\in\cA(S)$, 
 $$c(u,t)\leq \begin{cases} x(\delta^\uparrow(u))Z_u\leq 1/5  & \text{if $x(\delta^\uparrow(u))< \eps_F$}\\
  (1/2+9\eps_{1/2})(1-\eps_B) & \text{if 
 $x(\delta^\uparrow(u))>\eps_F$ and $u$ incident to bad edge}	\\
 1 + \eps_\eta & \text{otherwise, using \cref{lem:shared-edges}.}
 \end{cases}
 $$
 %(since $\eps_F \le 1/10$). Otherwise, if , then $x(\delta^\uparrow(u))\leq 1/2+9\eps_{1/2}$ (and $F_u=1-\eps_B$), so $c(u,t) \le (1/2+9\eps_{1/2})F_u$. If neither of these two cases hold, then, $c(u,t) \le 1 + \eps_\eta$ by \cref{lem:shared-edges}. Therefore,
Therefore, we can write,
 $$ c(\delta^\uparrow(W),t) \leq k(1/2+9\eps_{1/2})(1-\eps_B) + (|W|-k)(1+\eps_{\eta}).$$
 Now, to prove \eqref{eq:goalTdeltarigthup}, using \eqref{eq:xGEdeltalower}, it is enough to choose $\alpha$ and $\eps_B$ such that,
 $$
 (1+\alpha) \left(|W| - \eps_{\eta}/2 - k(1/2+\eps_{1/2})\right) \geq k(1/2+9\eps_{1/2})(1-\eps_B) + (|W|-k)(1+\eps_{\eta}),
 $$
 or equivalently,
 $$ |W|(\alpha-\eps_{\eta}) \geq k (\alpha/2 + 10\eps_{1/2} + \alpha \eps_{1/2} -\eps_B/2 - 9\eps_B \eps_{1/2} -\eps_{\eta}) + \frac{\eps_{\eta}}{2}(1+\alpha)$$
Since  every atom is adjacent to at most one bad edge, $k\leq |W|$
and $|W|\ge 1$, the inequality follows using
$\eps_B \ge  21 \eps_{1/2}$ and $\alpha > 2 \eps_\eta$ and $\eps_{1/2}\leq 0.0002$ and $\eps_\eta \leq \eps_{1/2}^2$.
\end{proof}

%\subsection{Bottom edges}
%
%\begin{lemma}
%For a polygon cut $S$, there is a matching from edges in $E^\rightarrow(S)$ to $\delta(S)$ such that in every near mincut $C$ in the polygon $S$ the edges in $C \cap \delta(S)$ are matched to at least capacity $x(C \cap \delta(S))$ in $E^\rightarrow(S)$. Furthermore, edges $e$ in $\delta^\rightarrow(S)$ are matched to at most fraction $(2+3\epsilon_{\eta})x_e$ for $\epsilon_\eta \le 1/10$. 
%\end{lemma}\label{lem:bottom-matching}
%\begin{proof}
%	Match every edge in $\delta^\rightarrow(S)$ to every edge in $\delta^\uparrow(S)$, and then multiply its capacity by $(2+3\epsilon_{\eta})$. By \cref{lem:shared-edges} every cut $C$ has 
%	$$x(C \cap \delta^\uparrow(S)) \le 1+\epsilon_\eta$$ 
%	Therefore since $\delta(C) \ge 2$, $C \cap \delta^\uparrow(S)$ is matched to at least fraction $1-\epsilon_\eta$ before multiplying its capacity. After multiplication
%	$$(1-\epsilon_\eta)(1+3\epsilon_\eta) \ge 1+\epsilon_\eta \ge x(C \cap \delta^\uparrow(S))$$
%\end{proof}
